\chapter{CS II常见问题}

\textbf{问题1：IF 阶段如何对 block RAM 进行信号输入。}

解答：因为 block RAM 只能进行同步读操作，为了可以在 IF 阶段获得需要的指令，取指的 PC 值必须提前一周期提供给 block RAM。因此输入给 block RAM 的 PC 值不能使用 PC 寄存器的输出，而要使用 PC 寄存器的输入。需要特别注意的是，PC 寄存器的输入分别是 PC 的初始值（寄存器初始化）、PC 本身（流水线阻塞）、PC 的递增值（普通指令）和跳转指令的目标地址（分支跳转或者跳转指令），因此 block RAM 的输入也需要是四种情况中的一种。

\textbf{问题2：寄存器的值在时钟上升沿到来前发生变化/寄存器的值提前一拍发生变化。}

解答：这往往是因为混用使用阻塞赋值 \texttt{=} 和非阻塞赋值 \texttt{=>} 导致的。非阻塞赋值会将 \texttt{reg} 综合为触发器，然后在时钟边沿修改触发器的值；但是阻塞赋值可能会将 \texttt{reg} 综合为线路，然后在当前时钟周期就修改线路的值，导致输出意外的提前一个周期向后传播。如果对阻塞赋值的机制不甚了解，请尽量避免使用阻塞赋值，也尽量避免混用阻塞赋值和非阻塞赋值。

\textbf{问题3：\texttt{kernel} 运行的时候遇到内存访问权限异常。}

解答：在 \texttt{kernel} 执行过程中遇到内存访问权限异常，往往是因为没有为 \texttt{kernel} 进程分配堆栈：\texttt{sp} 寄存器没有初始化，或者没有为堆栈预分配内存空间。

\textbf{问题4：处理器差分测试时没有输出。}

解答：该错误可能是处理器内部流水线阻塞导致的，往往是因为处理器内核和内存模块之间握手失败，导致处理器内核一直在等待内存模块的反馈。

\textbf{问题5：\texttt{project} 在差分测试时遇到错误，但是无法在该指令附近定位问题。}

解答：该错误可能是异常中断处理机制发生特权级切换时引入的。当时钟中断等异常中断事件发生时，正在执行的内存写操作、寄存器写回操作应该被无效化。如果在特权级切换的时钟周期没有立刻关闭内存读写使能和寄存器写使能，就会导致寄存器文件和内存中的值与差分测试模型不一致。这些操作往往在退出异常中断处理程序之后才会被差分测试发现，从而提高了错误定位的难度。

\textbf{问题6：下板之后遇到处理器卡死在 \texttt{PC=0x5c}。}

解答：处理器执行的程序中有一个用于内存初始化的循环，并对初始化结果进行正确性检查，如果不正确就会死循环在 0x5c 地址。因为 DDR MIG 模块对于时序要求很高，如果处理器设计不规范，会导致处理器时序问题严重，进而导致 DDR MIG 读写结果错误，从而在内存初始化检查环节陷入 0x5c 死循环。此时有如下几种可以尝试的解决方法：

\begin{itemize}
    \item 长按 reset 按钮，重新进行内存初始化。
    \item 修改处理器内存初始化程序，将内存校验部分代码删除，或者校验失败重新内存初始化。
    \item 修改处理器中不规范的写法，修复处理器的时序问题，特别是 critical warning。
    \item 使用 global 的方法综合 DDR MIG，提高处理器整体的时序稳定性。
    \item 使用 Vivado 的高级优化选项，提高处理器的时序稳定性。
\end{itemize}
